\paragraph{一族\trDoubleCover{}。}\label{sub:family-double-covers}
我们现在来构造一个\trLefschetz\trBifibration，它包含 $\tilde{\pi}\colon\widetilde{E}\to\widetilde{S}$ 连同其对其他\trBranchPoint——包括\trDegenerate（\trSingular）的\trBranchPoint——的类似物。设 $\scr{W}\subset\CC$ 是半径 $1-2\delta$ 的\trDisk。考虑 $\scr{W}\times S$ 沿对角线分歧的\trDoubleCover，以及 $\scr{W}\times E$ 沿 $\pi$ 的图像分歧的\trDoubleCover（或更精确地说，对角线／图像中有一个坐标落在 $\scr{W}\subset S$ 之内的那些部分）：
\begin{equation*}
\begin{gathered}
\widetilde{\scr{S}}=\{(w,y)\in\scr{W}\times\CC:\ y^2+w\in S\},\\
\widetilde{\scr{E}}=\{(w,y,x)\in\widetilde{\scr{S}}\times E:\ y^2+w=\pi(x)\}.
\end{gathered}
\end{equation*}
$\widetilde{\scr{S}}$ 自带一个自然的\trComplexStructure，而 $\widetilde{\scr{E}}$ 带有一个\trAlmostComplex{}结构；两种情形下，边界都是弱\trConvex{}的。我们有下列\trPseudoholomorphic{}映射
\begin{equation*}
\begin{array}{ll}
\left.\begin{aligned}
\iota_{\widetilde{\scr{S}}}(w,y)&=(w,-y)\colon\widetilde{\scr{S}}\to\widetilde{\scr{S}}\\
\iota_{\widetilde{\scr{E}}}(w,y,x)&=(w,-y,x)\colon\widetilde{\scr{E}}\to\widetilde{\scr{E}}
\end{aligned}\right\} & \text{\trCovering\trInvolution；}\\[14pt]
\left.\begin{aligned}
\beta_{\widetilde{\scr{S}}}(w,y)&=y^2+w\colon\widetilde{\scr{S}}\longrightarrow S\\
\beta_{\widetilde{\scr{E}}}(w,y,x)&=x\colon\widetilde{\scr{E}}\longrightarrow E
\end{aligned}\right\} & \text{\trQuotient{}映射，忘掉 $w$；}\\[14pt]
\varpi(w,y,x)=(w,y)\colon\widetilde{\scr{E}}\longrightarrow\widetilde{\scr{S}}, & \text{对每个 $w$，$\pi$ 的\trLift；}\\[14pt]
\left.\begin{aligned}
\psi(w,y)&=w\colon\widetilde{\scr{S}}\longrightarrow\scr{W}\\
\Psi(w,y,x)&=w\colon\widetilde{\scr{E}}\longrightarrow\scr{W}
\end{aligned}\right\} & \text{到参数 $w$ 的诸\trProjection。}
\end{array}
\end{equation*}
注意 $\Psi=\psi\circ\varpi$。$\varpi$ 的\trSingularPoint{}为
\begin{equation}\label{eq:family-crit-sets}
\mathrm{Crit}(\varpi)=\bigsqcup_{x\in\mathrm{Crit}(\pi)}\{y^2+w=\pi(x)\},\qquad
\mathrm{Critv}(\varpi)=\bigsqcup_{z\in\mathrm{Critv}(\pi)}\{y^2+w=z\}.
\end{equation}
在第一个集合之中，$\Psi$ 的\trCriticalPoint{}是形如 $(w=\pi(x),0,x)$ 的那些，这意味着 $\mathrm{Critv}(\Psi)=\mathrm{Crit}(\pi)$。如果我们在满足 $\pi(x)=z$ 的一点 $x\in\mathrm{Crit}(\pi)$ 附近选取局部\trHolomorphic{}\trMorse{}坐标，则 $\widetilde{\scr{E}}$ 的定义方程就变成
\begin{equation}\label{eq:local-morse-equation}
w=z-y^2+x_1^2+\cdots+x_{n+1}^2;
\end{equation}
消去 $w$ 变量之后，便有 $\varpi(x,y)=(z-y^2+x_1^2+\cdots+x_{n+1}^2,\,y)$。由此容易看出，$\varpi$ 与 $\Psi$ 都有\trGeneric{}的\trCriticalPoint。此外，$\Psi$ 在\trCriticalPoint{}$(z,0,x)$ 处的\trHessian，若限制到 $\ker(D\varpi)$——即切于 $y=0$ 的那个子空间——之上，仍然保持\trNondegenerate。这些计算表明，就\trAlmostComplex{}几何而论，
\begin{equation}\label{eq:family-bifibration-diagram}
\begin{tikzcd}[column sep=large]
\widetilde{\scr{E}} \ar[r,"\varpi"] \ar[rr,bend left=16,"{\Psi=\psi\circ\varpi}"]
 & \widetilde{\scr{S}} \ar[r,"\psi"] & \scr{W}
\end{tikzcd}
\end{equation}
满足\trLefschetz\trBifibration{}的一切性质（实际上，这是\trBifibration{}的一个特别简单的例子，因为那里没有假\trCriticalPoint）。尽管这整个讨论本质上是初等的，把要点以一种更直观的形式重新表述一遍也许仍值得一做。对每个 $w$，限制
\begin{equation}\label{eq:family-fibre-restriction}
\varpi_w\colon\widetilde{\scr{E}}_w=\Psi^{-1}(w)\longrightarrow\widetilde{\scr{S}}_w=\psi^{-1}(w)
\end{equation}
是 $\pi$ 沿\trDoubleBranchedCover{}$\beta_{\widetilde{\scr{S}}}|\widetilde{\scr{S}}_w\colon\widetilde{\scr{S}}_w\to S$ 的\trPullback（特别地，取 $w=*$，我们就回到了先前定义的 $\tilde{\pi}\colon\widetilde{E}\to\widetilde{S}$）\footnote{译注：原书此处印作“for $\pi=*$”；疑应为“$w=*$”（$\varpi_*$ 即先前定义的 $\tilde{\pi}$，而 $\pi=*$ 无意义，2x 放大实证）；此处径改。——译者注}。每当 $w\notin\mathrm{Critv}(\pi)$ 时，$\widetilde{\scr{E}}_w$ 是\trSmooth{}的，而 \eqref{eq:family-fibre-restriction} 有\trGeneric{}的\trCriticalPoint，它们成对出现，对应于 $\pi$ 的\trCriticalPoint。更精确地说，固定 $w$，则对每个 $x\in\mathrm{Crit}(\pi)$ 都对应有两点 $(w,\pm\sqrt{\pi(x)-w},x)$。当 $w$ 趋近\trCriticalValue{}$z=\pi(x)$ 时，这两个平方根相互靠拢，其相遇之处正是 $\widetilde{\scr{E}}_z$ 唯一的\trSingularPoint。
取使 $\scr{W}$ 成为一个带边\trExactSymplectic\trManifold{}的 $\theta_{\scr{W}}$、$\phi_{\scr{W}}$。我们将假定它们足够大。考虑把 $\theta_{\scr{W}}$ 的\trPullback（经 $\psi$）、$\kappa$ 的\trPullback（经第二个\trProjection）与 $\theta_S$ 的\trPullback（经 $\beta_{\widetilde{\scr{S}}}$）相加所得的 1-形式 $\theta_{\widetilde{\scr{S}}}$，并令 $\phi_{\widetilde{\scr{S}}}=d\theta_{\widetilde{\scr{S}}}$。我们已经看到，$\phi_{\widetilde{\scr{S}}}|_{\widetilde{\scr{S}}_*}=\phi_{\widetilde{S}}$ 是辛的\footnote{译注：原书此处印作“$\theta_{\scr{W}}|\widetilde{\scr{S}}_*=\phi_{\widetilde{S}}$”；疑应为“$\phi_{\widetilde{\scr{S}}}|_{\widetilde{\scr{S}}_*}=\phi_{\widetilde{S}}$”（依 \eqref{eq:shrunk-base-forms}，$\psi$ 项在\trFiber{}上为零，$\phi_{\widetilde{\scr{S}}}$ 限制在 $\widetilde{\scr{S}}_*$ 上恰为 $\phi_{\widetilde{S}}$；且“限制到其他任何\trFiber{}同理”只对 2-形式成立，4x 放大实证原书如此）；此处径改。——译者注}，对限制到其他任何\trFiber{}上亦有同样的结论。于是，由于我们已加上了从底空间\trPullback{}而来的大\trSymplecticForm{}$\phi_{\scr{W}}$，$\phi_{\widetilde{\scr{S}}}$ 本身是辛的。用与\cref{thm:transport-isotopy-realization}相同的论证（特别参见 \eqref{eq:liouville-inwards-criterion}），\trLiouville\trVectorField{}将沿位于 $\partial\scr{W}$ 之上的\trBoundaryFace{}$\partial^v\widetilde{\scr{S}}$ 指向内侧。为处理另一张面，考虑部分\trTrivialization
\begin{equation}\label{eq:partial-trivialization-family}
\scr{W}\times\widetilde{S}\supset\scr{W}\times\widetilde{U}\longrightarrow\widetilde{\scr{S}},\qquad
(w,y)\longmapsto\bigl(w,\sqrt{y^2+*-w}\,\bigr)
\end{equation}
它定义在子集 $\widetilde{U}=\{|y^2+*|\ge 1-\delta\}\subset\widetilde{S}$ 之上；我们选取当 $w=*$ 时给出恒同映射的那个唯一的平方根。\eqref{eq:partial-trivialization-family} 的像是\trFibrewise{}边界 $\partial^h\widetilde{\scr{S}}$ 的一个邻域。由于在 $\scr{W}\times\widetilde{U}$ 上有 $|\sqrt{y^2+*-w}\,|\ge\sqrt{\delta}$，这个部分\trTrivialization{}与给定的\trSymplecticStructure{}及 1-形式原函数\trCompatible，其含义是：$\theta_{\widetilde{\scr{S}}}$ 经 \eqref{eq:partial-trivialization-family} 的\trPullback{}是 $\theta_{\scr{W}}+(\theta_S|\widetilde{U})$。由此容易得出：$\widetilde{\scr{S}}$ 是一个\trCorneredBoundary{}的\trExactSymplectic\trManifold，而且 $\psi$ 的诸\trHorizontal\trTangentSpace{}平行于 $\partial^h\widetilde{\scr{S}}$。
相应地，令 $\theta_{\widetilde{\scr{E}}}=\varpi^*(\theta_{\scr{W}}+\kappa)+\beta_{\widetilde{\scr{E}}}^*\theta_E$。与前面相同的论证表明：若 $\phi_{\scr{W}}$ 足够大，则 $\phi_{\widetilde{\scr{E}}}=d\theta_{\widetilde{\scr{E}}}$ 是辛的，除非可能在 $\Psi$ 的\trCriticalPoint{}附近。如果现在像 \eqref{eq:local-morse-equation} 中那样在局部坐标 $(x,y)$ 下看一个\trCriticalPoint，那么 $\phi_{\widetilde{\scr{E}}}$ 就等于 $d\kappa$（在 $y$-坐标上）加上 $\phi_E$（在 $x$-坐标上）再加上 $\phi_{\scr{W}}$ 经 $\Psi(y,x)=z-y^2+x_1^2+\cdots+x_{n+1}^2$ 的\trPullback；这显然是一个\trKaehler{}形式。接下来我们考察 $\widetilde{\scr{E}}$ 的各张\trBoundaryFace。(i) 在由 $|w|=1-2\delta$ 定义的那张面——即 $\Psi^{-1}(\partial\scr{W})$——上，前面用于 $\partial^v\widetilde{\scr{S}}$ 的论证适用，并表明负\trLiouville\trVectorField{}指向内侧。(ii) 在面 $|y^2+w|=1$（即 $x\in\partial^vE$ 所在之处）的一个邻域内，有与 \eqref{eq:partial-trivialization-family} 类似的部分\trTrivialization。\trLiouville{}场的性态因而归结为 $\scr{W}\times\widetilde{E}$ 的情形。注意，这张面是 $\widetilde{\scr{E}}$ 相对于 $\Psi$ 而取的\trHorizontal{}边界的一部分；这个\trTrivialization{}还表明，$\Psi$ 的\trSymplectic\trConnection{}在它附近是平凡的。(iii) 剩下的那张面由那些 $x\in\partial^hE$ 的点 $(w,y,x)$ 组成。在这样的一个点处，考虑
\begin{equation}\label{eq:differential-splitting-family}
D\beta_{\widetilde{\scr{E}}}\colon T\widetilde{\scr{E}}=T\widetilde{\scr{E}}^h\oplus T\widetilde{\scr{E}}^v\longrightarrow TE=TE^h\oplus TE^v.
\end{equation}
第二个\trDirectSummand{}$T\widetilde{\scr{E}}^v$ 是沿 $\varpi$ 的\trFiber{}的\trTangentSpace，因而\trIsomorphic{}地映到 $TE^v$ 上。第一个\trDirectSummand{}$T\widetilde{\scr{E}}^h$ 是它关于 $\phi_{\widetilde{\scr{E}}}$ 的\trOrthogonal{}补；由于这两个形式之差是从 $\varpi$ 的底空间\trPullback{}而来，这其实与关于 $\beta_{\widetilde{\scr{E}}}^*\phi_E$ 的\trOrthogonal{}补相同。因此，$D\beta_{\widetilde{\scr{E}}}$ 保持 \eqref{eq:differential-splitting-family} 的\trSplitting。由于映射 $\beta_{\widetilde{\scr{E}}}$ 与 $\partial^hE$ \trTransverse，沿我们这张\trBoundaryFace{}的\trTangentSpace{}是 $T(\partial^hE)$ 在 \eqref{eq:differential-splitting-family} 之下的\trPreimage，从而包含整个 $T\widetilde{\scr{E}}^h$。这表明与 $\varpi$ 相伴的\trSymplectic\trConnection{}的\trHorizontal\trTangentSpace{}在 $(w,y,x)$ 处平行于边界；出于显然的理由，这对 $\Psi$ 也成立。此外，$T\widetilde{\scr{E}}$ 中的负\trLiouville{}向量与 $TE$ 上相应\trVectorField{}的差别只在于一个 $T\widetilde{\scr{E}}^h$ 分量，而由于该分量切于边界，\trLiouville{}向量必定指向内侧。我们现在已经证明：$\widetilde{\scr{E}}$ 是一个\trCorneredBoundary{}的\trExactSymplectic\trManifold，而 $\varpi$、$\Psi$ 都是\trExact{}的\trFibrationWithSingularities；这就验证了 \eqref{eq:family-bifibration-diagram} 满足\trLefschetz\trBifibration{}定义中的辛方面。为将来之用，我们记录一个关于 $\Psi$ 的几何的简单观察：
\begin{lemma}\label{thm:monodromy-involution-agreement}
设 $\lambda$ 是在 $\scr{W}$ 中绕半径为 $|*|=1-3\delta$ 的\trCircle{}正向一周的\trLoop。考虑 $\widetilde{\scr{E}}_*=\widetilde{E}$ 相伴的单值\trAutomorphism。在子集 $\widetilde{E}^{\mathrm{shrink}}$ 上，这一映射与\trInvolution{}$\iota_{\widetilde{E}}$ 相一致。
\end{lemma}
\begin{proof}
在任何满足 $|w|=1-3\delta$ 与 $|\pi(x)|\le 1-4\delta$ 的点 $(w,x,y)$ 处，我们有 $|y|=|\sqrt{\pi(x)-w}\,|\ge\sqrt{\delta}$，因此 $d\kappa$ 对\trSymplecticForm{}的贡献消失。于是，$\Psi$ 的\trSymplectic\trConnection{}的限制就等于平凡\trConnection{}经\trProjection{}到 $\scr{W}\times E$ 的\trPullback。因此，沿 $\lambda(t)=(1-3\delta)e^{2\pi it}$ 的\trParallelTransport{}就成为
\begin{equation}\label{eq:monodromy-parallel-transport}
h_{\lambda|[0;t]}\colon(*,y,x)\longmapsto\bigl(\lambda(t),\sqrt{y^2+\lambda(0)-\lambda(t)},x\bigr)=\Bigl(\lambda(t),\,y\sqrt{\frac{\pi(x)-\lambda(t)}{\pi(x)-\lambda(0)}},\,x\Bigr),
\end{equation}
其中平方根的选取使 \eqref{eq:monodromy-parallel-transport} 在 $t=0$ 时为恒同。当 $t$ 从 $0$ 走到 $1$ 时，$\lambda(t)-\pi(x)$ 绕原点描出一个\trLoop，故平方根变号。
\end{proof}
